\documentclass[11pt,a4paper]{article}

\usepackage[T2A]{fontenc}
\usepackage[utf8]{inputenc}
\usepackage[russian,english]{babel}
\usepackage{geometry}
\usepackage{amsmath,amssymb}
\usepackage{hyperref}
\usepackage{booktabs}
\usepackage{enumitem}

\geometry{margin=2.5cm}

\title{Angel-in-Pocket 2.0\\
Центральный агент принятия решений с интеграцией GRA-устойчивости}
\author{qqewq}
\date{\today}

\begin{document}
\selectlanguage{russian}
\maketitle

\begin{abstract}
Angel-in-Pocket 2.0 — это центральный AI-агент, который объединяет бизнес, науку, производство и творчество в единую систему принятия решений. В качестве бэкенда устойчивости используется GRA-Core, который анализирует сценарии, выявляет противоречия и обнуляет нестабильные ветви. Архитектура ориентирована на расширяемый монорепозиторий с модулями для ангела, бизнеса, финансов, подписок, науки, искусства, backend- и frontend-слоёв.
\end{abstract}

\noindent\textbf{Ключевые слова:} ИИ-агенты, бизнес-автоматизация, подписки, бухгалтерия, налоги, наука, искусство, устойчивость, GRA.

\section{Введение}

Angel-in-Pocket 2.0 представляет собой центрального агента, который принимает пользовательскую цель и превращает её в рабочий план действий. Пользователь задаёт намерение, ограничения и предпочтения, а система на их основе строит сценарии, оценивает их и выбирает устойчивый вариант.

Ключевая особенность проекта — интеграция GRA-Core как слоя устойчивости. Этот слой нужен для ранжирования стратегий, поиска противоречий и удаления неэффективных ветвей до того, как они начнут расходовать ресурсы.

\section{Постановка задачи}

Цель системы можно представить как преобразование множества входов в единую операционную политику:
\[
(V, P, C) \rightarrow B \rightarrow G \rightarrow \Pi,
\]
где $V$ — видение пользователя, $P$ — модель предпочтений, $C$ — ограничения, $B$ — бизнес-канва, $G$ — граф процессов, а $\Pi$ — итоговая политика исполнения.

Задача агента состоит не только в генерации ответа, но и в поддержании согласованности между бизнесом, финансами, наукой и творческими процессами.

\section{Архитектура репозитория}

Монорепозиторий разделён на несколько доменных слоёв:
\begin{itemize}[leftmargin=1.2cm]
    \item \textbf{angel} — центральное ядро агента и планировщик.
    \item \textbf{business} — арбитраж, продажи, маркетинг, стратегия и тактика.
    \item \textbf{finance} — бухгалтерия, учёт, доходы, расходы и налоги.
    \item \textbf{billing} — подписки, тарифы, платёжные провайдеры и удержание клиентов.
    \item \textbf{production} — планирование производства, закупки и себестоимость.
    \item \textbf{science} — анализ данных, гипотезы и экспериментальные сценарии.
    \item \textbf{arts} — музыка, изображение, текст и креативные задачи.
    \item \textbf{gra\_integration} — слой связи с GRA-Core.
    \item \textbf{backend} — API и серверная логика.
    \item \textbf{frontend} — пользовательский интерфейс и дашборды.
\end{itemize}

Каждый слой может работать независимо, но все важные решения проходят через общий слой устойчивости.

\section{Ядро агента}

Ядро ангела получает цель и контекст, после чего строит набор сценариев. Каждый сценарий оценивается по полезности, риску и устойчивости.

Формально полезность можно записать как:
\[
U(a \mid s) = \sum_{i=1}^{n} w_i \cdot \phi_i(a, s),
\]
где $a$ — действие, $s$ — состояние системы, $\phi_i$ — критерии оценки, а $w_i$ — веса приоритетов.

Выбранное действие определяется как:
\[
a^* = \arg\max_a U(a \mid s).
\]

\section{Бизнес-модуль}

Бизнес-слой отвечает за маркетинг, продажи, стратегию, тактику и арбитраж.  
В проекте предусмотрена логика типа «купить дешевле, продать дороже», а также анализ рынка, воронки продаж и продуктовые стратегии.

Стратегию можно описать как оптимизацию прибыли при ограничениях риска:
\[
\Pi = R - O - \tau,
\]
где $R$ — выручка, $O$ — операционные расходы, а $\tau$ — налоги.

\section{Финансы и бухгалтерия}

Финансовый слой ведёт учёт денежных потоков, журнал операций, базовый план счетов и простые налоговые сценарии.  
Бухгалтерия нужна не как отдельная формальность, а как часть общей логики принятия решений.

Для любого проекта можно рассматривать:
\begin{itemize}[leftmargin=1.2cm]
    \item доходы,
    \item расходы,
    \item обязательства,
    \item денежный остаток,
    \item налоговую нагрузку.
\end{itemize}

\section{Подписки и монетизация}

Billing-модуль поддерживает разные модели монетизации:
\begin{itemize}[leftmargin=1.2cm]
    \item подписка по тарифам,
    \item разовая оплата,
    \item usage-based billing,
    \item донаты,
    \item корпоративные планы,
    \item кастомные офферы.
\end{itemize}

Подписочная модель может быть записана как:
\[
T_i = \{p_i, b_i, r_i, c_i\},
\]
где $p_i$ — цена, $b_i$ — биллинговый период, $r_i$ — удержание, а $c_i$ — риск оттока.

\section{Наука}

Научный модуль предназначен для работы с данными, гипотезами и моделями.  
Он может применять symbolic regression, анализировать таблицы, сравнивать гипотезы и формировать отчёты.

Одна из типовых задач — поиск формулы по данным:
\[
\min_f \; \mathcal{L}(f; D) + \lambda \cdot \Omega(f),
\]
где $\mathcal{L}$ — ошибка на данных, а $\Omega$ — штраф за сложность модели.

\section{Искусство}

Творческий модуль поддерживает генерацию текста, музыки и визуального контента.  
Он может работать как соавтор, куратор или редактор, а не только как генератор.

Для креативных решений полезно учитывать новизну:
\[
U_{\text{creative}}(a \mid s) = \log P(a \mid s) + \beta \cdot N(a, s),
\]
где $N(a,s)$ — мера новизны, а $\beta$ задаёт баланс между качеством и оригинальностью.

\section{GRA-устойчивость}

GRA-Core в этой архитектуре выполняет роль слоя проверки и обнуления.  
Если сценарий слишком противоречив, дорог или нестабилен, он должен быть пересмотрен или удалён.

Условие сохранения ветви можно записать так:
\[
\text{Keep}(x) \iff U(x) - K(x) > 0,
\]
где $U(x)$ — полезность, а $K(x)$ — стоимость, риск и противоречие.

\section{Заключение}

Angel-in-Pocket 2.0 — это не просто ассистент и не просто бизнес-бот.  
Это единый агент принятия решений, который соединяет бизнес, финансы, науку, искусство и производство под управлением GRA-устойчивости.

Архитектура проекта рассчитана на расширение и позволяет добавлять новые домены без разрушения ядра.

\begin{thebibliography}{9}

\bibitem{latex}
LaTeX Project. \emph{LaTeX: A document preparation system}. \url{https://www.latex-project.org/}

\bibitem{ctan}
CTAN. \emph{The Comprehensive TeX Archive Network}. \url{https://ctan.org/}

\bibitem{gra}
GRA-Core: Unified Hierarchical Stability Library. Репозиторий и документация проекта.

\end{thebibliography}

\end{document}